%!TEX root = ../../springer-mono-format/book.tex
%!TEX root = ../../thesis-trier.tex


ere, in \cref{sync_max:thm:IS_SCC} (resp.~\cref{sync_max:thm:IS_SCC})
we show that strong connectivity (resp.~strong $\frac{3}{2}$-connectivity)
of $I(S)$ implies complete reachability and distinguishability
of pairs of $2$-sets (resp.~complete distinguishability) in a transformation semigroup $S$.
Furthermore, we conclude in \cref{thm:primitive} that strong connectivity
implies absence of any non-trivial congruence,
which gives minimality in terms of automata as stated in \cref{cor:primitive}.

\begin{theorem}
	\label{sync_max:thm:IS_SCC}
	Let $M \le_M \mathcal T_n$ be a transformation monoid. If $I(M)$ is strongly connected, then
	$M$ is completely reachable and 
	every $2$-set $\{a,b\}$ is distinguishable
	from every $B \subseteq [n]$
	with $\{a,b\} \not\subseteq B$.	
%	two subsets $A, B \subseteq [n]$ 
%    with $|A|, |B| \ge 2$ and $A \not\subseteq B$ and $B \not\subseteq A$
%    are distinguishable. 
    In particular, two distinct $2$-subsets in $[n]$ are distinguishable.   
\end{theorem}
\begin{proof}
    Suppose $I(M)$ is strongly connected.
	Let $\mathcal A$ be an automaton with $\TM{\mathcal A} = M$
	constructed as in \cref{prelim:remarks_trans_semigrp_semiautomata}
	by selecting a generating set of~$M$ as a set of letters.
	Then, as $\Gamma_1(\mathcal A) = \Gamma_1(M)$ and $I(M)$ is a subgraph of $\Gamma_1(M)$ by \cref{sync_max:prop:IS_in_Gamma1} , 
	by \cref{sync_max:thm:char_compl_reach}, $M$ is completely reachable.
		
    Next, we can assume $|B| \ge 2$
    and, without loss of generality, suppose $a \notin B$.
	Let $x_0, x_1, x_2, \ldots, x_k$ be a shortest path in $I(M)$
	with $x_0 = a$ and $x_k = b$.
	By definition, there exist simple idempotent mappings $t_i \in M$
	with $\applyfunc{x_{i-1}}{t_i} = x_i$ for $i \in \{1,\ldots,k\}$.
% nach def von path alle verschieden
%	As the path is shortest possible, it contains no loops 
%	and so $|\{x_0, x_1, \ldots, x_k\}| = k + 1$.
	We distinguish two cases.
	
	\begin{enumerate}
	\item $B \not\subseteq \{x_0, x_1, \ldots, x_k\}$.
		
		Let $b_1, b_2 \in B$ be distinct with $b_1 \notin \{x_0, x_1, \ldots, x_k\}$
		and $t = \concat{t_1}{\concat{\cdots t_2}{t_k}}$.
		By \cref{prelim:eqn:concat-SI}
		we find $\applyfunc{\{a,b\}}{t} = \{b\}$
		and \[
		\applyfunc{\{b_1, b_2\}}{t} = \begin{cases}
			\{b_1, b\}   & \mbox{if } b_2 \in    \{x_0,x_1, \ldots, x_k\}, \\
			\{b_1, b_2\} & \mbox{if } b_2 \notin \{x_0,x_1, \ldots, x_k\}.
		\end{cases}
	    \]
		
	\item $B \subseteq \{x_0, x_1, \ldots, x_k\}$. 
		
		As $x_0 \notin B$,
		we have $B \subseteq \{x_1, \ldots, x_k\}$.
		Set $\overline t = \concat{t_2}{\concat{\cdots t_{3}}{t_k}}$.
		By \cref{prelim:eqn:concat-SI} we find
		$\applyfunc{B}{\overline t} = \{x_k\}$
		and $\{x_0, x_k\}  = \applyfunc{\{x_0, x_k\}}{\overline t}$ where $\{x_0, x_k\} = \{a,b\}$.
	\end{enumerate}
    In both cases, there exists a transformation distinguishing $\{a,b\}$ from $B$.
\end{proof}

\begin{remark} 
	To deduce complete reachability in \cref{sync_max:thm:IS_SCC}
	we must start from a transformation monoid and not merely from a transformation semigroup
	since otherwise the whole domain might not be reachable.
	For example, consider the semigroup of all transformations of defect at least one on the set~$[n]$.
	In this semigroup $[n]$ is not reachable,
	but since it contains all simple idempotents, $I(S)$
	equals the digraph where all pairs of vertices are connected by an edge 
	(but which implies that every proper and non-empty set is reachable).
\end{remark}

\begin{example}
	Let $t_i : [3] \to [3]$ be three distinct constant transformations
	with $t_i(0) = t_i(1) = t_i(2) = i$
	tor $i \in [3]$. The set $S = \{ t_1, t_2, t_3 \}$ forms a transformation semigroup.
	The digraph on~$[3]$ obtained by connecting point $a$ to $b$
	if and only if there exists an idempotent mapping $a$ to $b$
	is strongly $\frac{3}{2}$-connected.
	However, no pair of distinct non-empty and non-singleton 
	subsets of states is distinguishable 
	in  $S$ and $S$ is not completely reachable.
	In this sense, restricting ourselves to simple idempotents in the definition of $I(S)$
	is necessary.			
\end{example}

\Cref{sync_max:thm:IS_SCC} with \cref{sync_max:thm:lower_bound_SynA}
and the fact that $I(\mathcal A)$
is a subgraph of $I(\TM{\mathcal A})$ for SI-automata $\mathcal A$
yields the following corollary.

\begin{corollary}
\label{cor:IA_SCC_implies_CR}
	Let $\mathcal A$ be an $n$-state automaton.
	If $I(\TM{\mathcal A})$ (or $I(\mathcal A)$ if $\mathcal A$ is an SI-automaton) is strongly connected,
	then $\mathcal A$ is completely reachable and a minimal automaton
	for $\Syn(\mathcal A)$ has at least $\binom{n}{2} + 1$ states.
\end{corollary}

\begin{example}
Consider the strongly connected SI-automaton from \cref{sync_max:fig:SCC_not_sufficient}. 
Here, $G(\mathcal A) = I(\mathcal A) = I(\TM{\mathcal A})$.
So $I(\mathcal A)$ is strongly connected, but $\mathcal A$
is not completely distinguishable.		
\end{example}

Next, we state a sufficient condition for complete reachability
and complete distinguishability.


\begin{theorem}
	\label{sync_max:thm:IS_32SCC}
	Let $M \le_S \mathcal T_n$ be a transformation monoid. If $I(M)$ is strongly~$\frac{3}{2}$-connected and $n \ge 3$,
	then $M$ is completely reachable and completely distinguishable.
\end{theorem} 
\begin{proof}
	As strong $\frac{3}{2}$-connectivity implies strong connectivity,
	by \cref{sync_max:thm:IS_SCC}
	complete reachability follows.
	Next, let $A \subseteq [n]$ be a set containing
	two distinct states $p, q \in A$
	and $r \in [n] \setminus A$.
	By assumption there exists a path $q_0, q_1, \ldots, q_m$
	in $I(M)$ such that $r \notin \{q_0, q_1, \ldots, q_m\}$
	and $\{q_0, q_m\} = \{p,q\}$.
	Without loss of generality,
	suppose $q_0 = p$ and $q_m = q$.
	By the definition of $I(M)$ there exist
	simple idempotents $t_i \colon [n]\to [n]$
	for each $i \in \{1,\ldots,m\}$
	such that $\applyfunc{q_{i-1}}{t_i} = q_i$.
	Set $t = \concat{t_1}{\concat{\cdots t_2}{t_k}}$.  
	By \cref{prelim:eqn:concat-SI}
	we find $\applyfunc{\{p,q\}}{t} = \{q\}$, $\applyfunc{A}{t} \subseteq A$,
	and $\applyfunc{r}{t} = r$.
	Hence $|\applyfunc{A}{t}| < |A|$
	and $\applyfunc{r}{t} \notin \applyfunc{A}{t}$.
	From~\crefitem{sync_max:thm:char_distinguishable}{dist:item:compress}
	it follows that~$M$ is completely distinguishable.	
	\begin{comment}
		Now, assume $|S| + |T| > 4$. % two dsitnct
		Without loss of generality, there exists $a \in S \setminus T$. \todo{menge $X$ statt $[n]$?}
		Furthermore, let $b,c,d \in [n] \setminus \{a\}$ be three pairwise distinct elements such that
		$\{a,b\} \subseteq S$
		and $\{c,d\} \subseteq T$.
		
		If $I(S)$ is strongly $\frac{3}{2}$-connected,
		then there exists
		a path from $c$ to $d$ or from $d$ to $c$
		such that $a$ does not occur along the path.
		Without loss of generality, suppose $x_0, x_1, \ldots, x_k$ is a path in $I(S)$
		with $x_0 = c$, $x_k = d$ and $a \notin \{x_0, x_1, \ldots, x_k\}$.	
		Furthermore, by deleting any loops along
		the path, we can assume the elements $x_0, x_1, \ldots, x_k$ are pairwise distinct.
		Set $f = \concat{f_1}{\concat{\cdots f_2}{f_k}}$.
		%Then $|\applyfunc{\{c,d\}}{f}| = 1$.
		With Lemma~\ref{sync_max:lem:concat_SI}
		we deduce $\applyfunc{T}{f} = T \setminus \{x_0, x_1, \ldots, x_{k-1}\}$
		and $\applyfunc{S}{f} = S \setminus \{x_0, x_1, \ldots, x_{k-1}\}$.
		As $|\applyfunc{T}{f}| < |T|$
		we have $|\applyfunc{S}{f}| + |\applyfunc{T}{f}| < |S| + |T|$.
		Furthermore, we have $a \in \applyfunc{S}{f}$ as $a \notin \{x_0, \ldots, x_k\}$
		and $a \notin \applyfunc{T}{f}$.
		So $\applyfunc{S}{f} \ne \applyfunc{T}{f}$.
		
		Lastly, note that if $b \notin \{x_0, x_1, \ldots, x_k\}$,
		then $\{a,b\} \subseteq \applyfunc{S}{f}$.
		If  $b \in \{x_0, x_1, \ldots, x_k\}$,
		then $\applyfunc{b}{f} = d$
		and so $\{a,b\} \subseteq \applyfunc{S}{f}$. \todo{{a,d} drin?}
		Hence $3 \le |\applyfunc{S}{f}| + |\applyfunc{T}{f}|$.
		As $I(S)$ is completely reachable, it contains a constant mapping.
		So, we conclude that $I(S)$ has the DCSP property.
		
		
		Combining the DCSP property with the property that all pairs of $2$-sets
		are distinguishable, with Proposition~\ref{todo}
		we deduce that all pairs of non-singleton and non-empty subsets 
		of states are distinguishable in $I(S)$.	
	\end{comment}
\end{proof}

However, strong $\frac{3}{2}$-connectivity of $I(M)$ is not a necessary
condition for complete distinguishability, as shown by the next example.

\begin{example} \label{sync_max:ex:not-strongly-32-but-compl-dist}
	Consider the SI-automaton $\semiautomat{A}{\{1,2,3,4\}}{\Sigma}{\delta}$
	without permutational letters with automaton graph
	given in \cref{sync_max:fig:not-strongly-32-but-compl-dist}.
	This automaton graph 
	equals $I(\TS{\mathcal A})$,
	and it is not strongly $\frac{3}{2}$-connected
	since the state~$2$ is part of every path between $3$
	and any other state.
	However, the automaton is completely distinguishable.
	To see this, first, suppose $q \in A \setminus B$ for two non-empty and non-singleton subsets of states $A,B$
	for some state $q \ne 2$.
	Let~$u$ be a concatenation of the letters $a_{p,2}$
	with $p \in \{1,4,3\} \setminus \{q\}$
	in any order.
	Then $\delta(B, u) = \{2\}$
	and $\delta(A, u) = \{q,2\}$.
	Next, we have $A = \delta(Q, a_{4,1}) = Q \setminus \{4\}$
	and $B = \delta(Q \setminus \{2\}, a_{4,1}) = Q \setminus \{2,4\}$,
	and since the latter sets are distinguishable
	as shown before, we obtain
	that also $Q$ and $Q \setminus \{2\}$
	are distinguishable.
	
	\begin{figure}[ht]
		\centering
		\begin{tikzpicture}[node distance = 2.8em and 8em]
			\node[state] (1) {$1$};
			\node[state, right = of 1] (2) {$2$};
			\node[state, right = of 2] (3) {$3$};
			\node[state, above = of 2] (4) {$4$};
			
			\path[->] (1) edge [bend left=12,above] node {$a_{1,2}$} (2)
			(2) edge [bend left=12,above] node {$a_{2,3}$} (3)
			(3) edge [bend left=12,below] node {$a_{3,2}$} (2)
			(2) edge [bend left=12,below] node {$a_{2,1}$} (1);
			
			\path[->] (2) edge [bend left,left] node {$a_{2,4}$} (4)
			(4) edge [bend left,right] node {$a_{4,2}$} (2)
			(4) edge [bend right,above] node {$a_{4,1}$} (1);
		\end{tikzpicture}
		\caption[Pairs of $2$-sets distinguishable but not completely distinguishable]{A strongly connected and completely distinguishable SI-automaton $\mathcal A$ (self-loops omitted) where $I(\TS{\mathcal A})$ 
			is not strongly $\frac{3}{2}$-connected. See \cref{sync_max:ex:not-strongly-32-but-compl-dist}.}	
		\label{sync_max:fig:not-strongly-32-but-compl-dist}
	\end{figure}
\end{example}

Note that for $n \le 2$ every strongly connected graph is $\frac{3}{2}$-connected
and every transformation semigroup 
on $[n]$ is completely distinguishable (but possibly not completely reachable).

\begin{example}
	Strong $\frac{3}{2}$-connectivity of the $\Gamma_1$-graph
	is not sufficient for complete distinguishability.
	For example, consider the SI-automaton 
	from \cref{sync_max:fig:SCC_not_sufficient},
	in its $\Gamma_1$-graph all pairs of states
	are connected by an edge by \cref{sync_max:prop:Gamma1alledges}.
	So, it is  strongly $\frac{3}{2}$-connected.
	However, the automaton is not completely distinguishable.
\end{example}

Next, we show that strong connectivity of $I(S)$ also implies
that $S$ admits no non-trivial $S$-congruences.

\begin{theorem}\label{thm:primitive}
	Let $S \le_S \mathcal T_n$ be a transformation semigroup. If $I(S)$ is strongly connected,
	then $S$ admits no non-trivial $S$-congruence on $[n]$.
\end{theorem}
\begin{proof}
	Suppose $\sim$ is a non-trivial $S$-congruence on $[n]$.
	Pick some congruence class $A \subseteq [n]$, for which $A \ne [n]$,
	with $|A| \ge 2$. 
	As $I(S)$ is strongly connected, by \LEM~\ref{prelim:lem:partition_SCC}, there exists an edge $(p, q)$ in~$I(S)$
	with $p \in A$ and $q \notin A$. 
	So, we find a simple idempotent $t \in S$
	with $\applyfunc{p}{t} = q$.
	Let $r \in A \setminus \{p\}$.
	Then $r \sim p$ but $r = \applyfunc{r}{t} \not\sim \applyfunc{p}{t}$,
	contradicting the assumption that $\sim$
	is an $S$-congruence. So there exists no non-trivial $S$-congruence on~$[n]$.    
\end{proof}

From~\cref{thm:primitive} with~\cref{prelim:thm:char_minaut}
we get the following corollary.

\begin{corollary}\label{cor:primitive}
	Let $\semiautomat{A}{Q}{\Sigma}{\delta}$.
	If $I(\TS{\mathcal A})$ is strongly connected,
	then $\mathcal A_{q, F}$ is minimal for any $q \in Q$
	and non-empty and proper $F \subseteq Q$.
\end{corollary}
\begin{proof}
	First, if $I(\TS{\mathcal A})$ is strongly connected, 
	then $\mathcal A$ is completely reachable
	by~\cref{cor:IA_SCC_implies_CR} and so also strongly connected.
	Hence, $\mathcal A_{q,F}$ is accessible for any $q \in Q$
	and proper and non-empty $F \subseteq Q$.
	By~\cref{thm:primitive} 
	it admits no $\TS{\mathcal A}$-congruence
	so~\crefitem{prelim:thm:char_minaut}{thm:minaut:item:congruence} yields	that $\mathcal A_{q, F}$ is minimal.	
\end{proof}